% Agent class hierarchy for the RCC agent-based model
% Wrapped externally in \adjustbox{max width=0.85\textwidth}
\begin{tikzpicture}[
  >=Stealth,
  every node/.style={
    rectangle, rounded corners=3pt, draw, align=center,
    minimum height=0.65cm, inner sep=4pt, font=\small
  },
  abstract/.style={fill=gray!20, dashed, draw=gray!70},
  tcellnode/.style={fill=treatment!15, draw=treatment!50},
  tumornode/.style={fill=tumorcell!15, draw=tumorcell!50},
  innatenode/.style={fill=immunecell!15, draw=immunecell!50},
  stromalnode/.style={fill=accent!15, draw=accent!50},
  bus/.style={thick, gray!60},
  arr/.style={->, thick, gray!60},
  annot/.style={draw=none, fill=none, minimum height=0pt, inner sep=1pt,
                font=\footnotesize\itshape},
]

% ===== Row 0: Root =====
\node[abstract] (root) at (5, 0) {\texttt{repast4py.Agent}};

% ===== Row 1: Cell + SexHormone / Cytokine =====
\node[abstract]    (cell)     at (2, -1.8) {\texttt{Cell}};
\node[stromalnode] (hormone)  at (11, -1.8) {\texttt{SexHormone}};
\node[stromalnode] (cytokine) at (11, -3.2) {\texttt{Cytokine}};

% Root -> level-1 bus
\draw[bus] (root.south) -- (5, -0.8);
\draw[bus] (2, -0.8) -- (11, -0.8);
\draw[arr] (2, -0.8) -- (cell.north);
\draw[arr] (11, -0.8) -- (hormone.north);
\draw[arr, dashed] (11, -0.8) -- (11, -1.2) -- (cytokine.north);

\node[annot] at (cell.west)      [left=2pt]  {abstract};
\node[annot] at (hormone.east)   [right=2pt] {immediate moves};
\node[annot] at (cytokine.east)  [right=2pt] {placeholder};

% ===== Row 2: Cell's direct children (first tier) =====
\coordinate (cbus) at (2, -2.7);
\draw[bus] (cell.south) -- (cbus);
\draw[bus] (-2.5, -2.7) -- (8, -2.7);

\node[abstract]   (tcell) at (-2.5, -3.8) {\texttt{TCell}};
\node[tumornode]  (tumor) at (0.2,  -3.8) {\texttt{TumorCell}};
\node[innatenode] (treg)  at (2.5,  -3.8) {\texttt{TregCell}};
\node[innatenode] (dc)    at (5,    -3.8) {\texttt{DendriticCell}};
\node[innatenode] (pdc)   at (8,    -3.8) {\texttt{PlasmacitoidDC}};

\foreach \n in {tcell, tumor, treg, dc, pdc} {
  \draw[arr] (\n.north |- cbus) -- (\n.north);
}

\node[annot] at (tcell.south) [below=1pt] {hormone perception, TCR};
\node[annot, font=\footnotesize\color{immunecell}] at (dc.south) [below=1pt]
     {+\,PhagocyticMixin};
\node[annot, font=\footnotesize\color{immunecell}] at (pdc.south) [below=1pt]
     {+\,PhagocyticMixin};

% ===== Row 3: Cell's direct children (second tier) =====
\coordinate (bus2drop) at (5, -2.7);
\coordinate (bus2) at (5, -5.4);
\draw[bus] (bus2drop) -- (bus2);
\draw[bus] (0.2, -5.4) -- (13.2, -5.4);

\node[innatenode]  (m1)     at (0.2,  -6.5) {\texttt{MacrophageM1}};
\node[innatenode]  (m2)     at (2.5,  -6.5) {\texttt{MacrophageM2}};
\node[innatenode]  (nk)     at (4.8,  -6.5) {\texttt{NaturalKiller}};
\node[innatenode]  (mast)   at (6.8,  -6.5) {\texttt{MastCell}};
\node[innatenode]  (neutro) at (8.8,  -6.5) {\texttt{Neutrophil}};
\node[stromalnode] (adipo)  at (10.8, -6.5) {\texttt{Adipocyte}};
\node[stromalnode] (blood)  at (13.2, -6.5) {\texttt{Blood}};

\foreach \n in {m1, m2, nk, mast, neutro, adipo, blood} {
  \draw[arr] (\n.north |- bus2) -- (\n.north);
}

\node[annot, font=\footnotesize\color{immunecell}] at (m1.south) [below=1pt]
     {+\,PhagocyticMixin};
\node[annot] at (blood.south) [below=1pt] {has\_step\,=\,False};

% ===== Row 3 (left): T-cell subclasses =====
\coordinate (tbusy) at (-2.5, -7.2);
\draw[bus] (tcell.south) -- ++(0, -1.3);
\draw[bus] (-8, -7.2) -- (3, -7.2);

\node[tcellnode] (cd8n) at (-8,   -8.3) {\texttt{CD8NaiveTCell}};
\node[tcellnode] (ctc)  at (-5.3, -8.3) {\texttt{CytotoxicTCell}};
\node[tcellnode] (cd4n) at (-2.5, -8.3) {\texttt{CD4NaiveTCell}};
\node[tcellnode] (th1)  at (0.3,  -8.3) {\texttt{CD4Helper1TCell}};
\node[tcellnode] (th2)  at (3,    -8.3) {\texttt{CD4Helper2TCell}};

\foreach \n in {cd8n, ctc, cd4n, th1, th2} {
  \draw[arr] (\n.north |- tbusy) -- (\n.north);
}

\node[annot] at (ctc.south) [below=1pt] {CD8+};
\node[annot] at (th1.south) [below=1pt] {Th1};
\node[annot] at (th2.south) [below=1pt] {Th2};

% ===== Legend (bottom-right) =====
\begin{scope}[shift={(11.5, -7.5)}]
  \node[draw=none, fill=none, font=\small\bfseries, minimum height=0pt,
        inner sep=0pt, anchor=north west] at (-0.2, 0) {Legend};
  \node[abstract, font=\footnotesize, minimum width=2.2cm]
        at (1, -0.6) {Abstract};
  \node[tcellnode, font=\footnotesize, minimum width=2.2cm]
        at (1, -1.3) {T~cells\,(5)};
  \node[tumornode, font=\footnotesize, minimum width=2.2cm]
        at (1, -2.0) {Tumor\,(1)};
  \node[innatenode, font=\footnotesize, minimum width=2.2cm]
        at (1, -2.7) {Innate\,(7)};
  \node[stromalnode, font=\footnotesize, minimum width=2.2cm]
        at (1, -3.4) {Stromal / signal\,(4)};
  \draw[gray!50, rounded corners=4pt]
        (-0.4, 0.3) rectangle (2.5, -3.8);
\end{scope}

\end{tikzpicture}
